\documentclass[10pt]{article}

%%% Doc layout
\usepackage{fullpage} 
\usepackage{booktabs}       % professional-quality tables
\usepackage{microtype}      % microtypography
\usepackage{parskip}
\usepackage{times}
\usepackage{graphicx}

%% Hyperlinks always black, no weird boxes
\usepackage[hyphens]{url}
\usepackage[colorlinks=true,allcolors=black,pdfborder={0 0 0}]{hyperref}

%%% Math typesetting
\usepackage{amsmath,amssymb}

%%% Write out problem statements in blue, solutions in black
\usepackage{xcolor}
\newcommand{\officialdirections}[1]{{\color{purple} #1}}

%%% Avoid automatic section numbers (we'll provide our own)
\setcounter{secnumdepth}{0}

%% --------------
%% Header
%% --------------
\usepackage{fancyhdr}
\fancyhf{}
\setlength{\headheight}{15pt}
\fancyhead[C]{\ifnum\value{page}=1 Tufts CS 145 BDL - 2026s - HW4 Submission \else \fi}
\fancyfoot[C]{\thepage} % page number
\renewcommand\headrulewidth{0pt}
\pagestyle{fancy}

%% --------------
%% Begin Document
%% --------------
\begin{document}

~\\ %% add vertical space
{\Large{\bf Student Name: TODO}}

~\\ %% add vertical space
{\bf Collaboration Statement:}

Total hours spent: TODO

I consulted the following human, textbook, or AI resources:
\begin{itemize}
\item TODO
\item TODO
\item $\ldots$	
\end{itemize}

By turning this document in, I attest that I have followed the 
\href{https://www.cs.tufts.edu/cs/145/2026s/index.html#collaboration}{[CS 145 Collaboration Policy]}. All solutions represent my own work. No solution text in this document was directly provided by other humans or artificial agents. 

\tableofcontents

\newpage

\subsection{1a: Plot of BCE Reconstruction Loss over Iterations for plain AE}

\renewcommand{\figurename}{Fig.}
\renewcommand{\thefigure}{1a}
 \begin{figure}[!h]
     \centering
     \begin{tabular}{c c}
     arch = [32]
     & 
     arch = [512]
     \\
     \includegraphics[height=6cm,width=0.44\textwidth,keepaspectratio]{example-image-a.pdf}
     &
     \includegraphics[height=6cm,width=0.44\textwidth,keepaspectratio]{example-image-a.pdf}
     \end{tabular}
     \label{fig:1a}
\caption{TODO BRIEF CAPTION HERE (one line only).
}%endcaption
 \end{figure}


\subsection{1b: Plot of L1 Error in Reconstructions over Iterations for plain AE}

\renewcommand{\thefigure}{1b}
 \begin{figure}[!h]
     \centering
     \begin{tabular}{c c}
     arch = [32]
     & 
     arch = [512]
     \\
     \includegraphics[height=6cm,width=0.44\textwidth,keepaspectratio]{example-image-b.pdf}
     &
     \includegraphics[height=6cm,width=0.44\textwidth,keepaspectratio]{example-image-b.pdf}
     \end{tabular}
     \label{fig:1b}
\caption{TODO BRIEF CAPTION HERE (one line only).
}%endcaption
 \end{figure}

\newpage


\subsection{1c: Generated images using AE-trained decoder}

\renewcommand{\figurename}{Fig.}

\renewcommand{\thefigure}{1c}
 \begin{figure}[!h]
     \centering
     \begin{tabular}{c c}
     arch = [32]
     & 
     arch = [512]
     \\
     \includegraphics[height=6cm,width=0.44\textwidth,keepaspectratio]{example-image-c.pdf}
     &
     \includegraphics[height=6cm,width=0.44\textwidth,keepaspectratio]{example-image-c.pdf}
     \end{tabular}
     \label{fig:1c}
\caption{TODO BRIEF CAPTION HERE (one line only).
}%endcaption
\end{figure}
 

\subsection{1d: Code vectors in 2D space for AE-trained decoder}

\renewcommand{\thefigure}{1d}
 \begin{figure}[!h]
     \centering
     \begin{tabular}{c c}
     arch = [32]
     & 
     arch = [512]
     \\
     \includegraphics[height=6cm,width=0.44\textwidth,keepaspectratio]{example-image-c.pdf}
     &
     \includegraphics[height=6cm,width=0.44\textwidth,keepaspectratio]{example-image-c.pdf}
     \end{tabular}
     \label{fig:1d}
\caption{TODO BRIEF CAPTION HERE (one line only).
}%endcaption
\end{figure}

\subsection{1e: Short Answer: Assess Train vs. Test Performance of both architectures.}

\officialdirections{
Comment on the relative values of the metrics in Fig 1a and 1b between train and test sets. Are there signs of overfitting? Underfitting?
}

TODO YOUR ANSWER HERE.

\newpage


\subsection{2a: Plot of BCE Reconstruction Loss over Iterations for VAE}

\renewcommand{\figurename}{Fig.}
\renewcommand{\thefigure}{2a}
 \begin{figure}[!h]
     \centering
     \begin{tabular}{c c}
     arch = [32]
     & 
     arch = [512]
     \\
     \includegraphics[height=6cm,width=0.44\textwidth,keepaspectratio]{example-image-a.pdf}
     &
     \includegraphics[height=6cm,width=0.44\textwidth,keepaspectratio]{example-image-a.pdf}
     \end{tabular}
     \label{fig:2a}
\caption{TODO BRIEF CAPTION HERE (one line only).
}%endcaption
 \end{figure}


\subsection{2b: Plot of L1 Error in Reconstructions over Iterations for VAE}

\renewcommand{\thefigure}{2b}
 \begin{figure}[!h]
     \centering
     \begin{tabular}{c c}
     arch = [32]
     & 
     arch = [512]
     \\
     \includegraphics[height=6cm,width=0.44\textwidth,keepaspectratio]{example-image-b.pdf}
     &
     \includegraphics[height=6cm,width=0.44\textwidth,keepaspectratio]{example-image-b.pdf}
     \end{tabular}
     \label{fig:2b}
\caption{TODO BRIEF CAPTION HERE (one line only).
}%endcaption
 \end{figure}

\newpage


\subsection{2c: Generated images using VAE-trained decoder}

\renewcommand{\figurename}{Fig.}

\renewcommand{\thefigure}{2c}
 \begin{figure}[!h]
     \centering
     \begin{tabular}{c c}
     arch = [32]
     & 
     arch = [512]
     \\
     \includegraphics[height=6cm,width=0.44\textwidth,keepaspectratio]{example-image-c.pdf}
     &
     \includegraphics[height=6cm,width=0.44\textwidth,keepaspectratio]{example-image-c.pdf}
     \end{tabular}
     \label{fig:2c}
\caption{TODO BRIEF CAPTION HERE (one line only).
}%endcaption
\end{figure}
 

\subsection{2d: Code vectors in 2D space for VAE-trained decoder}

\renewcommand{\thefigure}{2d}
 \begin{figure}[!h]
     \centering
     \begin{tabular}{c c}
     arch = [32]
     & 
     arch = [512]
     \\
     \includegraphics[height=6cm,width=0.44\textwidth,keepaspectratio]{example-image-c.pdf}
     &
     \includegraphics[height=6cm,width=0.44\textwidth,keepaspectratio]{example-image-c.pdf}
     \end{tabular}
     \label{fig:2d}
\caption{TODO BRIEF CAPTION HERE (one line only).
}%endcaption
\end{figure}

\subsection{2e: Short Answer: Contrast 2D codes learned by AE and VAE.}

\officialdirections{
Comment on the differences in the learned encoding spaces shown in Fig 1d versus Fig 2d. Can you explain the key differences in terms of how each model treats the variable z?
}

TODO YOUR ANSWER HERE.

\end{document}
